% thm-norm-shape: limit-shapes.tex:341-354 (label thm:norm-shape)
\paragraph{Paper excerpt.}
\begin{verbatim}
\begin{theorem}[Limit shape for a norm]\label{thm:norm-shape}
Let $d\geq2$, let $\gauge$ be a norm on~$\R^d$, let $\xi(x)\in\partial \gauge(x)$ be chosen at the sites $x\in\Z^d\setminus\{0\}$, let $\drift>0$ satisfy~\eqref{eq:ellipticity}, and let $r_n$ be as in~\eqref{eq:radius-norm}. Then, almost surely, the centrally excited random walk with norm~$\gauge$ has the following properties.
\begin{enumerate}[label=\textup{(\roman*)}]
\item \underline{\emph{Shape}}: for every $0<\eta<1$ and all sufficiently large~$n$,
\begin{equation*}
\{x\in\Z^d:\gauge(x)<(1-\eta)r_n\}\subset A_n\subset\{x\in\Z^d:\gauge(x)<(1+\eta)r_n\}\, .
\end{equation*}
\item \underline{\emph{Local times}}: as $n\to\infty$,
\begin{equation*}
\frac1{r_n}\max_{x\in\Z^d}\bigl|\ell_n(x)-2d\drift(r_n-\gauge(x))_+\bigr|\longrightarrow0\, .
\end{equation*}
\item \underline{\emph{Recurrence}}: the walk visits every site of~$\Z^d$ infinitely often.
\end{enumerate}
\end{theorem}
\end{verbatim}
\paragraph{Lean statement.}
\begin{verbatim}
theorem CERW.Frozen.norm_shape {d : ℕ} (hd : 2 ≤ d)
    {Ψ : EuclideanSpace ℝ (Fin d) → ℝ} (hΨ : CERW.IsNorm Ψ)
    {ξ : Site d → EuclideanSpace ℝ (Fin d)}
    (hξ : ∀ x : Site d, x ≠ 0 → CERW.IsSubgradient Ψ (CERW.toSpace x) (ξ x)) (hξ0 : ξ 0 = 0)
    {ε : ℝ} (hε : 0 < ε) (hεΨ : ∀ i : Fin d, ε * Ψ (CERW.coordVec i) < 1 / (d : ℝ))
    {Ω : Type u} [MeasurableSpace Ω] (μ : Measure Ω) [IsProbabilityMeasure μ]
    (X : ℕ → Ω → Site d) (hX : CERW.IsDriftCERW μ ε ξ X) :
    let r : ℕ → ℝ := fun n =>
      ((d + 1) * n / (2 * d * ε * CERW.normBallVolume Ψ)) ^ ((1 : ℝ) / (d + 1))
    ∀ᵐ ω ∂μ,
      (∀ η : ℝ, 0 < η → η < 1 → ∀ᶠ n : ℕ in atTop,
        {x : Site d | Ψ (CERW.toSpace x) < (1 - η) * r n} ⊆ ↑(CERW.departureRange (X · ω) n) ∧
        (↑(CERW.departureRange (X · ω) n) : Set (Site d)) ⊆
          {x | Ψ (CERW.toSpace x) < (1 + η) * r n}) ∧
      (∀ η : ℝ, 0 < η → ∀ᶠ n : ℕ in atTop, ∀ x : Site d,
        |(CERW.localTime (X · ω) n x : ℝ) - 2 * d * ε * max (r n - Ψ (CERW.toSpace x)) 0|
          ≤ η * r n) ∧
      (∀ x : Site d, ∃ᶠ j in atTop, X j ω = x)
--
\end{verbatim}
\paragraph{Transcription.}
Paper macros: \verb!\drift! is $\kappa$ (Lean \verb!ε!), \verb!\gauge! is $\Psi$ (Lean \verb!Ψ!), \verb!\nf{a}{b}! is $a/b$. \emph{Object mapping.} \begin{itemize} \item $\R^d$ is \verb!EuclideanSpace ℝ (Fin d)!; $\Z^d$ is \verb!LatticeProb.Site d! ($=$ \verb!Fin d → ℤ!), embedded by \verb!CERW.toSpace!, so $\gauge(x)$ for a site $x$ is \verb!Ψ (CERW.toSpace x)!; $e_i$ is \verb!CERW.coordVec i!. The standing assumption $d\ge2$ (lines 85 and 104) is \verb!hd : 2 ≤ d!. \item $\gauge$ a norm: \verb!hΨ : CERW.IsNorm Ψ! (subadditive, absolutely homogeneous, zero only at $0$). Nonnegativity, evenness and Lipschitz continuity follow (Lipschitz is machine-checked), so $|B_\gauge|$ is finite and positive. \item Chosen subgradients: \verb!ξ : Site d → EuclideanSpace ℝ (Fin d)! with \verb!hξ : ∀ x, x ≠ 0 → CERW.IsSubgradient Ψ (CERW.toSpace x) (ξ x)!. \verb!IsSubgradient Ψ x ξ! is $\forall y,\
